\documentclass[../../main.tex]{subfiles}
\begin{document}

\paragraph{【新課程】}
平面上で，距離$a$の2定点$A$，$B$を始点とする単位ベクトル(大きさが1のベクトル)
$\overrightarrow{AP}$，$\overrightarrow{BQ}$が
それぞれ始点のまわりに同じ向きに回転運動をなし，
$\overrightarrow{AP}$，$\overrightarrow{BQ}$の$\overrightarrow{AB}$から測った回転角は
それぞれ$2\theta$，$\theta$である．
$\overrightarrow{PQ}$の大きさの最大値を$a$で表わせ．

\paragraph{【旧課程】}
二次方程式$x(x-a)+m(x-b)(x-c)=0$が$m$のすべての正の値に対して実根をもつための必要十分条件を
$a$，$b$，$c$の大小関係によって表わせ．
ただし$a$，$b$，$c$は相異なる正数とする．

\end{document}
